\chapter{Curvature, Energy, and Hypergraph Dynamics}
\label{ch:curvature-hypergraph}

\section{Overview}

With the semantic manifold $\mathcal{M}$ (Ch.~\ref{ch:semantic-manifold}),
RSVP fields $(\Phi,\mathbf{v},S)$ (Ch.~\ref{ch:rsvp-fields}),
cognitive rewrite operators (Ch.~\ref{ch:lsystem}),
and the semantic stack $\mathcal{S}$ (Ch.~\ref{ch:sheaf-hyperstructures})
in place, we now introduce the geometric and energetic structure required to
couple RSVP fields to Polyxan hypergraphs.

This chapter develops:
\begin{itemize}
  \item a notion of \emph{semantic curvature} on embedded hypergraphs,
  \item energy and tension tensors derived from RSVP fields,
  \item hypergraph-induced backreaction on curvature,
  \item the conservation laws that follow from diffeomorphism invariance,
  \item and the resulting \emph{hypergraph–field evolution equations}.
\end{itemize}

The primary goal is to define a geometric analogue of general relativity’s
stress–energy tensor for typed hypergraphs, but adapted to RSVP dynamics and
semantic flows. This will be the key to the full field–graph coupling in
Chapter~\ref{ch:coupling}.

\section{Embedded Hypergraph Geometry}

Let $H$ be a typed Polyxan hypergraph with vertex set $V$ and hyperedge set
$E$ (Ch.~\ref{ch:typed-hypergraphs}).  
A semantic embedding
\[
\iota_H : H \hookrightarrow \mathcal{M}
\]
assigns each node and hyperedge a region of $\mathcal{M}$.

\subsection{Node geometry}

Each node $v\in V$ is embedded as a point or a small region
$\iota_H(v)\subset \mathcal{M}$, with tangent structure inherited from
$T\mathcal{M}$.

The RSVP-induced metric $g_{\mu\nu}$ (Ch.~\ref{ch:rsvp-fields}) pulls back:
\[
g_{ab}^{(v)} 
:= 
\iota_H^\ast g_{\mu\nu}.
\]

\subsection{Hyperedge geometry}

A hyperedge $e=(v_1,\dots,v_k)\in E$ induces a geometric simplex-like object:
\[
\iota_H(e)
\subseteq 
\bigcup_i \iota_H(v_i)
\]
together with a higher-order incidence relation given by its type signature.

\subsection{Curvature of embedded hypergraphs}

Define the \emph{hypergraph curvature} at a node $v$ as:
\[
\mathcal{K}(v) 
:= 
\sum_{e\ni v}
\mathrm{Ric}\!\left(\iota_H(e)\right)
\]
where $\mathrm{Ric}$ is the Ricci form induced by $g$.

Similarly, curvature of an entire hypergraph is:
\[
\mathcal{K}[H]
=
\sum_{v\in V} \mathcal{K}(v).
\]

\section{Energy and Tension from RSVP Fields}

The RSVP Lagrangian (Ch.~\ref{ch:rsvp-fields}) determines field energy as:
\[
\mathcal{L}_{RSVP}
=
\frac{1}{2}(\nabla\Phi)^2
-
\frac{\beta}{2}S^{-1}(\nabla S)^2
+
\gamma(\nabla\Phi\cdot\mathbf{v})
+
V(\Phi,S).
\]

\subsection{RSVP stress tensor}

Define the RSVP stress–energy-like tensor:
\[
T_{\mu\nu}
=
\nabla_\mu\Phi\nabla_\nu\Phi
-
\frac{1}{2}g_{\mu\nu}(\nabla\Phi)^2
+
\beta S^{-1}
\left[
\nabla_\mu S \nabla_\nu S
-
\frac{1}{2}g_{\mu\nu}(\nabla S)^2
\right]
+
\gamma\,\nabla_{(\mu}\Phi\,v_{\nu)}.
\]

This tensor determines:
\begin{itemize}
  \item local smoothing forces,
  \item anisotropic tension,
  \item entropy exchange flows.
\end{itemize}

\section{Hypergraph Energy Functional}

To couple $H$ to RSVP dynamics we introduce a hypergraph energy functional.

\subsection{Node energy}

For each node $v$ define:
\[
E_v
=
\int_{\iota_H(v)}
T_{\mu\nu} g^{\mu\nu}\, d\mathrm{vol}.
\]

\subsection{Hyperedge tension}

Each hyperedge $e=(v_1,\dots,v_k)$ induces tension:
\[
\tau_e
=
\sum_{i<j}
d_g\!\left(\iota_H(v_i),\iota_H(v_j)\right),
\]
where $d_g$ is geodesic distance.

\subsection{Total hypergraph energy}

\[
E_H
=
\sum_{v\in V} E_v
+
\lambda\sum_{e\in E}\tau_e.
\]

$\lambda>0$ encodes the stiffness of hyperedge connections.

\section{Backreaction: Hypergraphs Affect Curvature}

The RSVP effective metric $g_{\mu\nu}$ is not primitive; it depends on the
fields and thus indirectly on the hypergraph through backreaction.

\subsection{Hypergraph-induced curvature}

Define the backreaction tensor:
\[
\Xi_{\mu\nu}
=
\frac{\delta E_H}{\delta g^{\mu\nu}}.
\]

This contributes to the curvature equation:
\[
\mathrm{Ric}_{\mu\nu}
-
\frac{1}{2}g_{\mu\nu}R
=
T_{\mu\nu}
+
\Xi_{\mu\nu}.
\]

This is the analogue of Einstein’s equation, but in RSVP form.

\section{Discrete Evolution Equations}

We now derive hypergraph evolution rules induced by RSVP curvature.

\subsection{Gradient flow of node embeddings}

Each node embedding evolves by:
\[
\partial_t \iota_H(v)
=
-\frac{\delta E_H}{\delta \iota_H(v)}.
\]

Explicitly:
\[
\partial_t \iota_H(v)
=
-
\left[
T_{\mu\nu} n^\nu
+
\lambda 
\sum_{e\ni v}
\operatorname{grad}_g\,d_g(\iota_H(v),\iota_H(v'))
\right].
\]

\subsection{Hyperedge relaxation}

Hyperedges evolve by smoothing in RSVP geometry:
\[
\partial_t \tau_e
=
-\frac{\delta E_H}{\delta \tau_e}
=
-
\lambda^{-1}
\tau_e.
\]

\section{Continuum Limit and Semantic Flux}

In the dense limit, nodes form a measure $\rho_X$ over $\mathcal{M}$.

\subsection{Continuum node density}

\[
\rho_X(x)
=
\sum_{v\in V}\delta(x-\iota_H(v)).
\]

\subsection{Flux equation}

Nodes move under RSVP flow $\mathbf{v}$ and cognitive updating:
\[
\partial_t \rho_X
+
\nabla\cdot(\rho_X \mathbf{v})
=
\sigma_{\mathrm{rewr}},
\]
where $\sigma_{\mathrm{rewr}}$ captures rewrite-induced mass creation and
annihilation.

\subsection{Semantic curvature flow}

Hypergraph tension induces curvature evolution:
\[
\partial_t g_{\mu\nu}
=
-2\left(
T_{\mu\nu}+\Xi_{\mu\nu}
\right).
\]

This is a generalization of Hamilton’s Ricci flow, but driven by semantic and
hypergraphic forces.

\section{Conservation Laws}

Diffeomorphism invariance of the action implies the following.

\subsection{Energy conservation}

\[
\nabla^\mu(T_{\mu\nu}+\Xi_{\mu\nu})=0.
\]

\subsection{Hypergraph momentum conservation}

\[
\nabla^\mu \Xi_{\mu\nu}=0
\quad\Rightarrow\quad
\text{nodes move geodesically except for rewrite forces}.
\]

\subsection{Entropy production}

RSVP entropy obeys:
\[
\partial_t S + \nabla\cdot(S\mathbf{v}) = \sigma_S,
\]
where $\sigma_S$ is hypergraph-induced semantic dissipation.

\section{Summary}

In this chapter we constructed:
\begin{itemize}
  \item the geometry of embedded Polyxan hypergraphs,
  \item the RSVP stress–energy tensor,
  \item hypergraph energy and tension functionals,
  \item curvature backreaction $\Xi_{\mu\nu}$,
  \item discrete node and hyperedge evolution equations,
  \item continuum flux laws,
  \item and conservation laws from diffeomorphism invariance.
\end{itemize}

These structures provide the energetic and geometric machinery needed for the
full RSVP–Polyxan coupling, which is the subject of the next chapter.
